\documentclass[11pt]{amsart}

\usepackage[T1]{fontenc}
\usepackage{lmodern}
\usepackage{microtype}
\usepackage{amsmath,amssymb,amsthm}
\usepackage{booktabs}
\usepackage[hidelinks]{hyperref}

\hypersetup{
  pdftitle={Order-Minimal Pendant-Free Counterexamples to Conjecture 2 of Nurdin et al.}
}

\newtheorem{theorem}{Theorem}
\newtheorem{proposition}[theorem]{Proposition}
\newtheorem{corollary}[theorem]{Corollary}

\newcommand{\tvs}{\operatorname{tvs}}
\newcommand{\wt}{\operatorname{wt}}

\title[Order-minimal pendant-free counterexamples]
{Order-Minimal Pendant-Free Counterexamples to Conjecture 2 of Nurdin et al.}

\date{}

\subjclass[2020]{05C78}
\keywords{total vertex irregularity strength, total labeling, counterexample,
graph atlas}

\begin{document}

\begin{abstract}
Let $B(G)$ be the degree-count lower bound conjectured by Nurdin et al.\ to
be equal to the total vertex irregularity strength $\tvs(G)$. We determine
all connected simple graphs with minimum degree at least two and order at
most seven for which $\tvs(G)>B(G)$. Exactly three occur, all on seven
vertices: $K_{3,4}$ and two explicitly given non-complete-multipartite
graphs. Each has $B(G)=2$ and $\tvs(G)=3$. The finite classification is
certified by explicit $B(G)$-labelings for the other $580$ atlas graphs,
while the three obstructions have elementary proofs. We also give a
regular-core join construction which, together with the seven-vertex
examples, yields non-complete-multipartite counterexamples of every order
at least seven. Since $K_{3,4}$ appears to be covered already by a 2005
complete-bipartite result, the contribution is the order-minimal
classification and the regular-core extension, not the first
counterexample.
\end{abstract}

\maketitle

\section{Results}

Let $G=(V,E)$ be a finite simple connected graph. A \emph{vertex-irregular
total $k$-labeling} is a map
\[
  \phi:V\cup E\longrightarrow\{1,\ldots,k\}
\]
for which the vertex weights
\[
  \wt_\phi(v)=\phi(v)+\sum_{uv\in E}\phi(uv)
\]
are pairwise distinct. The least such $k$ is the \emph{total vertex
irregularity strength} $\tvs(G)$, introduced by Ba\v{c}a et
al.~\cite{Baca2007}.

For each integer $j$, let $n_j$ be the number of vertices of degree $j$,
and let $\delta$ and $\Delta$ be the minimum and maximum degrees. Nurdin et
al.~\cite{Nurdin2010} proved the lower bound
\[
  B(G):=
  \max_{\delta\le i\le\Delta}
  \left\lceil
    \frac{\delta+\sum_{j=\delta}^{i}n_j}{i+1}
  \right\rceil
  \le \tvs(G)
\]
and conjectured equality in \cite[Conjecture~2]{Nurdin2010}. Susanto,
Simanjuntak, and Baskoro later produced infinite counterexample families
and asked whether a counterexample can have no pendant
vertices~\cite[Question~4.2]{Susanto2023}. That existence question may
already have been settled by the earlier complete-bipartite work of Wijaya
et al.~\cite{Wijaya2005}, which appears to contain $\tvs(K_{3,4})=3$,
whereas $B(K_{3,4})=2$; we have not been able to consult that paper, and
Susanto et al.\ do not cite it when posing Question~4.2. We therefore make
no priority claim for pendant-free existence.

Our first result gives the exact minimum-order classification.

\begin{theorem}\label{thm:classification}
Among connected simple graphs $G$ with $\delta(G)\ge2$ and $|V(G)|\le7$,
the inequality $\tvs(G)>B(G)$ holds for exactly three isomorphism classes.
They are $K_{3,4}$ and the graphs $G_1,G_2$ on $\{0,\ldots,6\}$ with
\begin{align*}
E(G_1)&=\{01,12,23,04,34,15,35,06,16\},\\
E(G_2)&=\{01,02,23,14,34,05,45,06,36\},
\end{align*}
where $ij$ denotes the edge $\{i,j\}$. For all three graphs,
\[
  B(G)=2
  \qquad\text{and}\qquad
  \tvs(G)=3.
\]
\end{theorem}

The graph6 strings of $G_1,G_2$, and the atlas representative of
$K_{3,4}$ are, respectively,
\[
  \texttt{FhdU?},\qquad
  \texttt{FpUK\_},\qquad
  \texttt{FreRW}.
\]
The graphs $G_1$ and $G_2$ are nonisomorphic: $G_1$ contains a triangle,
whereas $G_2$ is triangle-free. Neither is complete multipartite, since in
each graph $0\not\sim3$ and $3\not\sim1$, but $0\sim1$; nonadjacency is
transitive in every complete multipartite graph.

Our second result extends the endpoint-weight obstruction from complete
multipartite graphs to an arbitrary regular core.

\begin{theorem}\label{thm:family}
Let $t\ge1$ and $d\ge\max\{2t,t+2\}$. If $H$ is a
$(d-t-2)$-regular simple graph on $d$ vertices, then
\[
  G=H\vee\overline{K}_{2t+2}
\]
satisfies $B(G)=2<\tvs(G)$.
\end{theorem}

\begin{corollary}\label{cor:orders}
For every $n\ge7$, there is a connected non-complete-multipartite graph
$G$ of order $n$ and minimum degree at least two such that
$B(G)=2<\tvs(G)$.
\end{corollary}

\section{The three order-minimal graphs}

The degree sequence of each of $G_1$ and $G_2$ is
$(2,2,2,2,3,3,4)$. Hence
\[
 B(G_i)=\max\left\{
 \left\lceil\frac{2+4}{3}\right\rceil,
 \left\lceil\frac{2+6}{4}\right\rceil,
 \left\lceil\frac{2+7}{5}\right\rceil
 \right\}=2.
\]

\begin{proposition}\label{prop:small}
Neither $G_1$ nor $G_2$ admits a vertex-irregular total $2$-labeling.
\end{proposition}

\begin{proof}
Suppose that $G_i$ has such a labeling. Its four degree-$2$ vertices have
four distinct weights in $\{3,4,5,6\}$, so they use all four values. The
two degree-$3$ vertices must then have weights $7$ and $8$, and the unique
degree-$4$ vertex $c$ has weight $9$ or $10$.

Let $z$ be the degree-$2$ vertex of weight $3$. The label of $z$ and both
incident edge labels are $1$. In $G_2$, every degree-$2$ vertex is adjacent
to $c$ and to one degree-$3$ vertex. In $G_1$, exactly one degree-$2$
vertex is adjacent to both degree-$3$ vertices; it cannot be $z$, because
the degree-$3$ vertex of weight $8$ and all its incident edges have label
$2$. Thus, in either graph, $z$ is adjacent to $c$ and to a degree-$3$
vertex $a$.

The edge $cz$ has label $1$, so $c$ cannot have weight $10$. Hence
$\wt(c)=9$, and every contribution at $c$ other than $cz$ is $2$.
Similarly, $\wt(a)=7$, and every contribution at $a$ other than $az$ is
$2$. The other degree-$3$ vertex has weight $8$, so all four of its
contributions are $2$. From the displayed edge sets, every degree-$2$
vertex other than $z$ has both incident edges among these forced
label-$2$ edges. Its weight is therefore at least $1+2+2=5$, contradicting
the required degree-$2$ weight $4$.
\end{proof}

For $K_{3,4}$,
\[
 B(K_{3,4})=
 \max\left\{
 \left\lceil\frac{3+4}{4}\right\rceil,
 \left\lceil\frac{3+7}{5}\right\rceil
 \right\}=2.
\]
A total $2$-labeling would give seven distinct weights in the
seven-integer interval $[4,10]$, so the weights would exhaust that
interval. Weight $4$ forces a degree-$3$ vertex and its three incident
edges to have label $1$, while weight $10$ forces a degree-$4$ vertex and
its four incident edges to have label $2$. Those two vertices are
adjacent, a contradiction.

For completeness, the following certificates give total $3$-labelings.
Vertices are ordered $0,1,\ldots,6$; edge labels are read in the displayed
edge order for $G_1,G_2$, and in the order
\[
  03,04,05,06,13,14,15,16,23,24,25,26
\]
for $K_{3,4}$ with bipartition
$\{0,1,2\}\cup\{3,4,5,6\}$.

\begin{center}
\small
\begin{tabular}{@{}ccll@{}}
\toprule
Graph & vertex labels & edge labels & vertex weights \\
\midrule
$G_1$     & $3313233$ & $111111123$    & $7,9,3,6,4,5,8$ \\
$G_2$     & $3123333$ & $111111132$    & $9,3,4,7,6,5,8$ \\
$K_{3,4}$ & $2221222$ & $111211231233$ & $7,9,11,4,6,8,10$ \\
\bottomrule
\end{tabular}
\end{center}

Thus all three graphs have total vertex irregularity strength $3$.

\begin{proof}[Proof of Theorem~\ref{thm:classification}]
The Read--Wilson graph atlas~\cite{ReadWilson1998} contains one
representative of every simple graph of order at most seven. Filtering for
connected graphs of minimum degree at least two gives
\[
\begin{array}{c|rrrrr}
|V(G)|&3&4&5&6&7\\
\hline
\text{graphs}&1&3&11&61&507.
\end{array}
\]
Each of the other $580$ graphs admits an explicit vertex-irregular total
$B(G)$-labeling, located by exhaustive search as follows. Writing
$\wt_\phi(v)=1+\deg(v)+a_v+b_v$ with $a_v=\phi(v)-1\in[0,k-1]$ and
$b_v=\sum_{e\ni v}(\phi(e)-1)$, one enumerates the $k^{|E|}$ vectors of
edge offsets $\phi(e)-1$, which determine the $b_v$, and tests by greedy
interval matching on the length-$k$ intervals
$[\deg(v)+b_v,\deg(v)+b_v+k-1]$ whether the $a_v$ can be chosen to make
the weights pairwise distinct. Run over all $583$ graphs of the census at
$k=B(G)$, this search returns a labeling for every graph except
\texttt{FhdU?}, \texttt{FpUK\_}, and \texttt{FreRW}, and terminates in
well under a second; it needs no ancillary data. The preceding hand
proofs establish infeasibility at $k=2$ for those three graphs. Hence the
classification is complete.
\end{proof}

\section{A regular-core construction}

\begin{proof}[Proof of Theorem~\ref{thm:family}]
The $2t+2$ vertices in the independent set have degree $d$, while each
vertex of $H$ has degree
\[
  (d-t-2)+(2t+2)=d+t.
\]
Thus $G$ has order $n=d+2t+2$, minimum degree $d$, and maximum degree
$d+t$. For $d\le i<d+t$, the cumulative degree count is $2t+2$, and
\[
  1<\frac{d+2t+2}{i+1}
  \le\frac{d+2t+2}{d+1}\le2.
\]
At $i=d+t$, all vertices have been counted and
\[
  \frac{d+n}{d+t+1}
  =\frac{2d+2t+2}{d+t+1}=2.
\]
Therefore $B(G)=2$.

Assume that a vertex-irregular total $2$-labeling exists. The two degree
classes have weight intervals
\[
  [d+1,2d+2]
  \qquad\text{and}\qquad
  [d+t+1,2d+2t+2].
\]
These intervals overlap, since $d+t+1\le2d+2$, and their union is the
interval $[d+1,2d+2t+2]$. It contains exactly
\[
  (2d+2t+2)-(d+1)+1=d+2t+2=n
\]
integers. The $n$ distinct vertex weights must therefore exhaust this
interval. Its minimum forces a degree-$d$ vertex and all its incident
edges to have label $1$; its maximum forces a degree-$(d+t)$ vertex and
all its incident edges to have label $2$. These vertices lie on opposite
sides of the join and are adjacent, so their common edge would have both
labels. This contradiction proves $\tvs(G)>2$.
\end{proof}

\begin{proof}[Proof of Corollary~\ref{cor:orders}]
For $n=7$, take $G_1$. For $n\ge8$, set $t=1$, $d=n-4$, and
$H=\overline{C_d}$. Then $H$ is $(d-3)$-regular, so
Theorem~\ref{thm:family} applies. Moreover,
\[
  \overline{H\vee\overline{K}_4}=C_d\cup K_4.
\]
Since $d\ge4$, this complement is not a disjoint union of cliques; hence
the resulting graph is not complete multipartite.
\end{proof}

\begin{thebibliography}{99}

\bibitem{Baca2007}
M. Ba\v{c}a, S. Jendro\v{l}, M. Miller, and J. Ryan,
\emph{On irregular total labellings},
Discrete Math. \textbf{307} (2007), 1378--1388.
\href{https://doi.org/10.1016/j.disc.2005.11.075}
{doi:10.1016/j.disc.2005.11.075}.

\bibitem{Nurdin2010}
Nurdin, E. T. Baskoro, A. N. M. Salman, and N. N. Gaos,
\emph{On the total vertex irregularity strength of trees},
Discrete Math. \textbf{310} (2010), 3043--3048.
\href{https://doi.org/10.1016/j.disc.2010.06.041}
{doi:10.1016/j.disc.2010.06.041}.

\bibitem{ReadWilson1998}
R. C. Read and R. J. Wilson,
\emph{An Atlas of Graphs},
Clarendon Press, Oxford, 1998.

\bibitem{Susanto2023}
F. Susanto, R. Simanjuntak, and E. T. Baskoro,
\emph{Counterexamples to the total vertex irregularity strength's
conjectures},
Discrete Math. Lett. \textbf{12} (2023), 159--165.
\href{https://doi.org/10.47443/dml.2023.111}
{doi:10.47443/dml.2023.111}.

\bibitem{Wijaya2005}
K. Wijaya, Slamin, Surahmat, and S. Jendro\v{l},
\emph{Total vertex irregular labeling of complete bipartite graphs},
J. Combin. Math. Combin. Comput. \textbf{55} (2005), 129--136.

\end{thebibliography}

\end{document}